\section{Building Intuition}
\label{sec:v10-controlled-checks}
\label{sec:v10-markov-check}

A synthetic Markov task fixes the oracle and a known sufficient
state, so the budget question can be asked directly: when can
local node labels stand in for some root labels under a fixed
token budget?

A document is a length-\(128\) token sequence. At each position, a
hidden Markov chain selects one of \(K\) registers; the selected
register emits an observed token from a disjoint synonym palette.
The learner sees the token row, not the hidden registers. The
oracle \(f^\ast\) counts adjacent hidden-register changes.

\begin{table}[t]
\centering
\small
\begin{tabular}{@{}lcc@{}}
\toprule
 & Recoverable & Structural \\
\midrule
Hidden registers \(K\) & \(4\) & \(12\) \\
Document length & \multicolumn{2}{c}{\(128\) tokens} \\
Synonym tokens per register & \multicolumn{2}{c}{\(4\), disjoint palettes} \\
Switch probability \(h\) & \(0.039\) & \(0.079\) \\
Expected changepoints & \multicolumn{2}{c}{\({\sim}5\) and \({\sim}10\)} \\
Training documents & \multicolumn{2}{c}{\(10{,}240\)} \\
\bottomrule
\end{tabular}
\caption{\textbf{Markov data-generating process.} The hidden register
stays fixed with probability \(1-h\), otherwise switches to a uniformly
chosen alternative. The observed token is drawn uniformly from that
register's palette.}
\label{tab:v10-markov-dgp}
\end{table}

The sufficient state for any span is
\[
  (\mathrm{count},\mathrm{first},\mathrm{last}),
\]
where the count records transitions inside the span and the boundary
fields let parent merges add the one transition that may cross the
child boundary:
\[
(c_L,a_L,b_L)\otimes(c_R,a_R,b_R)
=
(c_L+c_R+\One\{b_L\ne a_R\},a_L,b_R).
\]
The state is ordered and associative. A scalar child count alone
cannot recover the boundary correction; the triple state can.

\begin{figure}[t]
\centering
\includegraphics[width=\linewidth]{markov_changepoint_combined.pdf}
\caption{\textbf{Counting register changes through a merge tree.}
A simplified \(12\)-token document with the same structure used
in the experiments at scale (\(128\) tokens). Hidden colors run
\(A,A,A\mid B,B,C\mid C,D,D\mid A,A,B\); the learner sees only the
word row beneath. Each leaf \(L_i\) holds a triple
\((\text{count}, \text{first}, \text{last})\). Merges add child
counts plus the boundary indicator
\(\One\{\text{last}_{\text{left}}\ne\text{first}_{\text{right}}\}\):
\(S_{12}\) adds the \(A\!\to\!B\) boundary (count
\(=0+1+1=2\)); \(S_{34}\) adds the \(D\!\to\!A\) boundary
(count \(=1+1+1=3\)); the root adds the \(C\!\to\!C\) boundary
(no change), giving count \(=5\). A scalar leaf score cannot
carry the boundary information that turns three local counts
into a correct root.}
\label{fig:v10-markov-walkthrough}
\end{figure}

Figure~\ref{fig:v10-markov-walkthrough} reads the mergeable
summary on a \(12\)-token version of the same data-generating
process. The full experiments use \(128\)-token documents with the
same hidden Markov mechanism; the small version makes every leaf
state, boundary correction, and merge visible. The node label that
supervises learning is the triple, not the scalar count: a learner
trained on scalar leaf counts cannot recover the boundary
information that the merge needs.

The completed allocation grid asks the budget question directly.
It uses leaf sizes \(128,64,32,16,8\), root-label shares from
\(100\%\) down to \(10\%\), and three local-allocation policies
that decide which tree nodes spend the reallocated mass. A
full-document label on a \(128\)-token sequence costs \(128\) token
observations; a label on a \(16\)-token leaf costs \(16\), so eight
such leaf labels spend the same full-document-equivalent mass as
one root label.

\begin{figure}[t]
\centering
\includegraphics[width=\linewidth]{markov_leaf_size_fixed_recoverable.pdf}
\caption{\textbf{Performance across leaf size and root-label budget.}
Tree-vs-oracle root MAE at each leaf size (x-axis) and root-label
budget share (y-axis); lower is better. The top row (R100) is the
all-root reference, with the full root-label budget and no local
labels. Reducing the root-label share down to R20 keeps MAE low
across most leaf sizes when the missing root-label mass is
backfilled with local labels. Performance only collapses in the
R10 row at small leaves, where the budget cannot support either
channel. Appendix~\ref{app:v10-markov} reports the matching
structural panel (\(K=12\) registers).}
\label{fig:v10-markov-budget-grid}
\end{figure}

Figure~\ref{fig:v10-markov-budget-grid} reads the resulting
surface: across a broad interior band of (leaf, root-share) cells,
reallocating missing root-label mass to local count labels keeps
root MAE near the all-root reference.
Appendix~\ref{app:v10-markov} gives the full leaf-size grid, the
three allocation policies, and provenance.

When the local state is known and node labels measure the state
needed for composition, node supervision can substitute for some
root supervision at fixed budget. The next section runs the same
trade on real party platforms, where the sufficient state is no
longer known in advance.
